% Additive interfaces for the 1272-coordinate source.
% These use the original exponent-24 constant H_*, not H_19.
% They do not alter any existing source equation or polynomial artifact.

\begin{lemma}[The original scalar certificate at order thirty-two]
\label{v5:scalar32-bridge}
Let $K=\Q(\sqrt2)$, $d=3-\sqrt2$, $D=d^2$, and
$E:y^2=x^3-Dx$, with the point $P_*$ of
Section~\ref{sec:degree8-original}. Put
\[
 \beta=X^2+2,\qquad
 r=\beta(\beta-1)(\beta-2401)(32\beta+343)(2\beta-1),
 \qquad X\in\Q.
\]
For $h_0,a\in\Z\setminus\{0\}$ and points $J_0=h_0P_*$ and $S=aP_*$,
take the
two rational coefficients, in the basis $1,\sqrt2$, of
$f(J_0),f(S),z(S)/z(J_0)$ in that order. Denote them by
$X_1,\ldots,X_6$, and put
\[
 (m_1,\ldots,m_6)=(1,1,1,1,32,32),\qquad l_i=X_i/r^{m_i}.
\]
Use the same ordered norm tree as the original scalar construction:
\[
 \begin{split}
 N_r(u,v)&=u^2+rv^2,\\
 n_0&=N_r(1,l_1),\quad n_1=N_r(l_2,l_3),\quad
 n_3=N_r(l_5,l_6),\\
 W&=N_r\bigl(N_r(n_0,n_1),N_r(l_4,n_3)\bigr).
 \end{split}
\]
Let
\[
 c\in\left\{\pm\prod_{p\in I}p:
                   I\subseteq\{2,3,5,7,11,13,17\}\right\}.
\]
If $b=crW^2\ne0$ satisfies the original pole predicate $\Psi(b)$,
then
\[
 M_5(\beta)^{32}\mid a,\qquad
 h(\beta)\le H_*|a|^{3/4}.
\]
Here $M_5$ and the exact constant $H_*$ are those of
Theorem~\ref{hyp:o8-height}; in particular
$1\le H_*<2^{4783227}$.

For every integral $X$, the scalar certificate is complete with
successively chosen positive indices $h_0,a$. One may require $h_0$
to be divisible by any prescribed positive integer and $a$ to be
arbitrarily divisible and arbitrarily large. After $S$ has been chosen,
the normalized-leaf conditions and the nonvanishing of $W$ persist when
$h_0$ is replaced by every integer in one suitable congruence class
$h\equiv h_0\pmod L$. The modulus $L$ may also be required to be
divisible by any prescribed positive integer.
\end{lemma}

\begin{proof}
The value $r$ is nonzero: $\beta\ge2$, and $\beta=2401$ would
give $X^2=2399$, which has no rational solution. At a prime
$p>76176$, the multiplier $c$ is a unit. Every prime in $M_5(\beta)$
has odd valuation in $r$, by the contact computation in the proof of
Theorem~\ref{hyp:o8-height}. One can also observe that a negative
valuation of $\beta=X^2+2$ is even, so only its finite contacts
contribute here.

Write $R=v_p(r)$. The pole obstruction in Lemma~\ref{hyp:o8-pole}
and $b\ne0$ make the odd valuation $R+2v_p(W)$ positive. In a
norm node $u^2+rv^2$ the two nonzero summands have opposite valuation
parity, so there is no cancellation. The constant leaf on the
all-left path first excludes $R<0$. Then $v_p(W)>-R/2$.
Descending the tree, the implication
\[
 v_p(N_r(u,v))>-tR,\quad t<1
 \quad\Longrightarrow\quad
 v_p(u)>-tR/2,\quad v_p(v)>-(t+1)R/2
\]
shows that every leaf $l_i$ has valuation greater than $-R$.
Thus $v_p(X_i)>(m_i-1)R\ge m_i-1$, and integrality of the
valuation gives $v_p(X_i)\ge m_i$. This is the original norm
argument with the two assigned orders changed to thirty-two.

The order-one tests put $J_0$ and $S$ in the formal group at each
place above $p$. Such a completion is unramified, and the formal
logarithm is an isometry on that group. Since
$[h_0]S=[a]J_0$, it follows that
\[
 v_{\mathfrak p}\!\left(\frac{z(S)}{z(J_0)}\right)
       =v_p(a/h_0).
\]
Both rational coefficients of this ratio have valuation at least
thirty-two, so $v_p(a/h_0)\ge32$. As $h_0$ is a nonzero integer,
$v_p(a)\ge32$. This proves $M_5(\beta)^{32}\mid a$.
Theorem~\ref{hyp:o8-height} now gives
$h(\beta)\le H_*M_5(\beta)^{24}\le H_*|a|^{3/4}$.
No exponent-nineteen height estimate is used.

For completeness, let $X\in\Z$. The integer $r$ is nonzero and
$|r|>1$. The original seven-prime squareclass basis, proved in the
finite-correction argument following
Lemma~\ref{lem:original-seven-parity-completeness}, permits a choice
of $c$ such that $cr>0$ and $cr$ is square in
$\Q_2,\Q_3,\Q_{17}$. Let $\mathcal S$ be the finite prime support
of the integer $cr$.

At the places above $\mathcal S$, choose $h_0$ sufficiently
divisible that $J_0$ is in arbitrarily deep formal subgroups.
This remains possible after requiring any fixed divisibility of
$h_0$, because these subgroups have finite index in the local
elliptic groups. Having fixed this nonzero $J_0$, choose $a$
sufficiently divisible that $S=aP_*$ is still deeper. The
coefficients of $f(J_0),f(S)$ and $z(S)/z(J_0)$ then attain every
finite lower valuation bound needed here. This includes ramified
and bad primes, by working in sufficiently deep formal subgroups.
We arrange that all normalized leaves are integral when $p\mid r$
and divisible by $p$ when $p\mid c$ but $p\nmid r$.
The index $a$ may simultaneously exceed any fixed bound.

At $p\mid r$, the root of the displayed tree is congruent to one
modulo $p$, because $r\equiv0$ and the all-left leaf is one.
At $p\mid c$, $p\nmid r$, all other leaves vanish modulo $p$,
so the root is again congruent to one. Thus $W$ is a unit on
$\mathcal S$, and in particular is nonzero. At these primes
$v_p(crW^2)=v_p(cr)\ge0$; outside $\mathcal S$ its valuation is
even. Hence $b$ has square reduced denominator, is positive, and
has the same local squareclasses as $cr$ at two, three, and
seventeen. Lemma~\ref{lem:original-seven-parity-completeness}
supplies the scalar pole witnesses.

Finally, with $S$ fixed, these scalar tests are rational functions
of $J_0$ defined at the chosen point. Their coefficient maps are
continuous in the product of the completions above each prime in
$\mathcal S$. The finite lower valuation bounds persist in
sufficiently small neighborhoods. A common finite-index subgroup
of the integer orbit gives a congruence $h\equiv h_0\pmod L$
contained in all those neighborhoods. Enlarging $L$ preserves
that property. The same unit calculation proves the asserted
stability for the final shared point.
\end{proof}

\begin{lemma}[The fixed multiplier seventeen with preparation zero]
\label{v5:preparation-zero-bridge}
Let $b\in\Q^\times$, put $B=17b$, and define $\Psi_0(b)$ by
the existence of $e,p,q,s\in\Q$ satisfying
\[
 p^2=B(1-e)(1+3e),\qquad h=e(B-b^2),\qquad
 Bq^2=(s^2-Bh)(s^2+3Bh),
\]
where $h$ is the displayed expression, not a further free choice.
Then $\Psi_0(b)$ excludes negative odd valuations of $b$ at every
prime outside $\{2,3,17\}$. Conversely, it holds whenever $b$ has
square reduced denominator, is positive and square in
$\Q_2,\Q_3,\Q_{17}$, and
\[
 b<17,\qquad v_2(b)\ge3,\qquad v_3(b)\ge1,
 \qquad v_{17}(b)\ge2.
\]
No additional nonzero condition on $e,p,q,s,h$ is imposed.
\end{lemma}

\begin{proof}
Taking the preparation witnesses $t=0$ and $U=b$ in
\eqref{eq:o8-pole} gives exactly the displayed equations.
The obstruction therefore follows from
Lemma~\ref{lem:original-seven-parity-soundness}, including its
degenerate witness cases.

For completeness put
\[
 V=B(B-b^2)=B^2(1-b/17).
\]
The real inequality makes $V$ positive. The valuation bounds make
$1-b/17$ a square at two, three, and seventeen: its difference
from one has valuation at least three at two and at least one
at the two odd primes. Thus $V$ is square at these three places.

Let $\mu$ and $\theta$ be the positive squarefree representatives
of $B$ and $V$. The prescribed squareclasses of $b$ give
$\mu\equiv17\pmod{24}$; in particular $\mu>1$ and its support
contains seventeen but neither two nor three. Also
$\theta\equiv1\pmod{24}$, and seventeen is absent from its
support. At a prime $\ell\mid\mu$ outside the three exceptions,
$v_\ell(b)$ is positive and odd, since the reduced denominator
of $b$ is square. Hence $1-b/17$ is a square unit there, so $V$
and $\theta$ are local squares. The same assertion at seventeen
was already arranged. These facts make the supports disjoint.

Conversely, if $\ell\mid\theta$, then
$\ell\notin\{2,3,17\}$. If $v_\ell(b)>0$, the valuation
of $V$ is $2v_\ell(b)$; if $v_\ell(b)<0$, it is
$3v_\ell(b)$, again even because the denominator is square.
Thus an odd valuation of $V$ forces $v_\ell(b)=0$ and
$1-b/17\equiv0\pmod\ell$. It follows that
$B=17b\equiv17^2\pmod\ell$ is a nonzero square unit.
This proves the other mutual-squareness condition.

Lemma~\ref{lem:original-seven-simultaneous-realization} now gives
the two rational points in the squareclass of $B$. To make the
conversion explicit, write $V=\theta c^2$ and scale the second
point to the curve with parameter $nV$. If the first abscissa is
$Ba^2$, scale both points and parameters by $a$ so that their
abscissas are $B$ and $Bs^2$, with $s\ne0$, and their curve
parameters are $e_0$ and $e_0V$. For their ordinates $y_1,y_2$, set
\[
 e=e_0/B,\qquad p=y_1/B,\qquad q=y_2/(B^2s),\qquad
 h=e(B-b^2).
\]
The first curve equation gives $p^2=B(1-e)(1+3e)$.
Since $e_0V=B^2h$, the second gives
$Bq^2=(s^2-Bh)(s^2+3Bh)$. These are the required witnesses.
The only arithmetic existence theorem used in this specialization
is the existing simultaneous-realization lemma, with its stated
restricted-converse and descent dependencies.
\end{proof}

\begin{lemma}[Real bounds independent of later progression moduli]
\label{v5:real-progression-bridge}
Let $G$ be the closure of the integer orbit
\[
 \{([n]P_*,[n]P_*^\sigma):n\in\Z\}
       \subset E(\mathbb R)\times E^\sigma(\mathbb R),
\]
and let $G^0$ be its identity component. Fix a finite collection of
rational functions formed from nonzero integral multiples and fixed
translates of either elliptic coordinate, with all denominators
nonzero as rational functions on the relevant curve. Suppose their
pole set is contained in a finite union of inverse images of finite
sets under these nonzero multiplication-and-translation maps on
either factor. There is a bound $C$ such that every nonconstant
integer progression contains arbitrarily large positive indices at
which all these functions are defined and their real absolute values
are at most $C$. The bound is chosen from the functions before the
progression modulus is specified.

If $m=|G/G^0|$, $m\mid L$, and $m\mid h_0$, the progression
$h\equiv h_0\pmod L$ is dense in $G^0$, and its positive tail
meets every nonempty open subset of $G^0$ infinitely often.
\end{lemma}

\begin{proof}
The real elliptic groups are compact Lie groups. Their closed
subgroup $G$ is a compact abelian Lie group with finitely many
components, and $G^0$ is a torus. The projection of $G^0$ onto
each elliptic identity circle is surjective. Indeed the projected
integer orbit is infinite, since $P_*$ has infinite order and each
real embedding is injective on the group of algebraic points.
If the connected projection were trivial, the projected $G$ would
be finite; a nontrivial connected closed subgroup of a circle is
the whole circle.

Each nonzero multiplication map on a circle has finite fibers.
Consequently each of the prescribed inverse images of a finite
set has Haar measure zero in every component of $G$. The finite
union of pole sets cannot cover a component. Choose in each
component a nonempty open neighborhood whose compact closure
avoids the poles. Continuity bounds all the fixed rational
functions on these finitely many compact sets by one constant $C$.

For a progression $h=r+kn$, its closure is
$[r]P_*+[k]G$. Multiplication by the positive integer $k$ maps
$G^0$ onto itself, so this closure contains an entire component
of $G$. The positive tail of the progression is dense in its
closure: compactness gives unbounded positive multiples tending
to the identity, so each negative multiple is also a limit of
positive multiples. It therefore meets the chosen neighborhood
in that component at arbitrarily large positive indices.

If $m\mid k$, then $[k]G=G^0$. If also $m\mid r$, the coset
is $G^0$ itself. Apply this with $k=L$ and $r=h_0$ for the final
assertion.
\end{proof}
